\section{Conclusion}

We presented a pure-inertial framework that estimates camera displacement
from continuous 200-Hz IMU streams and audits the complete measurement chain
used to train and test it. The analysis separates the 3-s target into an
acceleration integral and an unobservable initial-velocity term. More than
more than forty estimator variants occupied the same broad error plateau; regressing out the
low-frequency velocity proxy reduced mean validation error by 16\%, whereas
longer context and post-hoc session fusion did not recover it. A confirmed
pause supplies a zero-velocity boundary condition, but the present experiment
contains only one sub-2-s interval and its benefit deteriorates rapidly as
attitude and bias errors accumulate. The measured pause rate of 0.5--2.3 per
minute is therefore an acquisition limitation in the current data.

PhysNet combines analytic anchor-frame pre-integration, dense displacement
supervision, a stillness-conditioned velocity gate, lever-arm estimation and calibrated
uncertainty. It reduced synthetic-test error by 24\% relative to adapted
IMUNet with one seventh of the parameters and achieved the lowest
leave-one-session-out error. On the single held-out recording the two methods
were statistically tied; their blend reduced error from 49.7~mm for zero
motion to 42.2~mm ($R^2=0.279$). The next validation step is therefore
experimental: acquire several independent, prospectively sealed recordings with controlled
pauses, improve the attitude solution, and compare the chessboard reference
with an independent tracker before making a reference-grade trajectory claim.
